\documentclass[11pt,a5paper]{article}
\usepackage[utf8]{inputenc}
\usepackage[T1]{fontenc}
\usepackage{lmodern}
\usepackage[french]{babel}
\usepackage[margin=1.50cm]{geometry}
\usepackage{booktabs}
\usepackage{array}
\usepackage{xcolor}
\usepackage{tikz}
\usetikzlibrary{arrows.meta,positioning}
\usepackage[colorlinks=true,allcolors=blue]{hyperref}

%% pdflatex histoiresecretelivret.tex && pdflatex histoiresecretelivret.tex && rm *.toc *.out *.aux *.log

% Symboles d'enseignes (symboles mathematiques standard)
\newcommand{\carreau}{\ensuremath{\diamondsuit}}
\newcommand{\coeur}{\ensuremath{\heartsuit}}
\newcommand{\pique}{\ensuremath{\spadesuit}}
\newcommand{\trefle}{\ensuremath{\clubsuit}}

\title{\textbf{L'Histoire secrète et le tarot}\\
\large Enseignes, correspondances et discordances~\newline synthèse documentée}
\author{Gabriel Chandesris}
\date{\today}

\begin{document}

\maketitle

\begin{center}
\small Série de Jean-Pierre Pécau (Delcourt, 2005--2019), \\ 
dérivée d'\emph{Arcanes} (1998) et d'\emph{Arcane majeur} (2003).\\
Sources croisées : Wikipédia, miroir fr-academic, BDThèque, Top Comics,\\
Babelio (citations du tome 1), gaudry.be, blog \emph{Les Secrets du Tarot}.
\end{center}

\vspace{0.5em}

\section{Le système des Archontes}

Vers 3000 av.~n.~è., le chaman Père Loup remet à quatre de ses petits-enfants
quatre \og ivoires \fg{}, des cartes magiques taillées dans l'ivoire, avec une
règle absolue : un seul ivoire par archonte, jamais réunis. Les quatre
deviennent immortels et traversent l'Histoire à la tête de leurs
\og maisons \fg{} (Coupes, Lances, Épées, Deniers), en maniant la
\og rétro-synchronicité \fg{} --- concept inspiré de la synchronicité de
Carl Gustav Jung --- qui permet d'entrevoir l'avenir et de l'infléchir vers
des probabilités plus avantageuses.

Chaque archonte est défini par un triple ancrage :
\begin{itemize}
\item une \textbf{enseigne latine / tarot} (l'ivoire, la maison) ;
\item une \textbf{couleur française moderne} (\carreau{}, \coeur{}, \pique{}, \trefle{}) ;
\item un \textbf{atout majeur} du tarot de Marseille.
\end{itemize}

Dans le tome 1 (\emph{Genèse}), le discours du chaman désigne les ivoires sous
leurs formes les plus anciennes : Coupe, \textbf{Lance}, Épée et
\textbf{Bouclier} (\og Erlin, à toi l'ivoire du bouclier pour défendre les
trois autres \fg{}). Le vocabulaire glisse ensuite selon les époques
traversées : le bouclier chamanique devient le \textbf{Denier} des ères
monétaires ; la Lance est tantôt \og Bâton \fg{} (nom tarotique traditionnel
de la maison) tantôt \og Lance \fg{} (forme guerrière).

\section{Correspondances officielles de la série}

\begin{table}[h]
\centering
\begin{tabular}{ p{0.15\textwidth} p{0.15\textwidth} p{0.15\textwidth} p{0.27\textwidth} p{0.25\textwidth}}
\toprule
\textbf{Archonte} 	& \textbf{Couleur} 		& \textbf{Enseigne\newline latine} & \textbf{Atout~\newline majeur} & \textbf{Maison} \\
\midrule
Dyo 				& \coeur{} Cœur 		& Coupe 				& Le Pape (V) 			& Coupes \\
Reka 				& \pique{} Pique 		& Bâton / Lance 		& La Papesse (II) 		& Lances / Bâtons \\
Aker$^{*}$ 			& \carreau{} Carreau 	& Épée 					& L'Impératrice (III) 	& Épées \\
Erlin 				& \trefle{} Trèfle 		& Denier / Bouclier 	& L'Empereur (IV) 		& Deniers \\
\midrule
Guillaume de Lecce 	& (Joker) 				& --- 					& Le Mat~\newline (non numéroté) 	& 5\ieme{} maison : moines rouges \\
\bottomrule
\end{tabular}
\caption{Les cinq archontes et leurs triplets enseigne / couleur / atout.\\
$^{*}$ Orthographe \og Arker \fg{} chez certains blogs ; Wikipédia écrit \og Aker \fg{}.}
\end{table}

Remarque de structure : les quatre atouts originels sont exactement les
\textbf{arcanes majeurs II à V} du tarot de Marseille (La Papesse,
L'Impératrice, L'Empereur, Le Pape), soit le quatuor des figures d'autorité
temporelle et spirituelle. Le cinquième archonte, Guillaume de Lecce
(introduit au tome 3, modelé sur Guillaume III de Sicile), incarne le
\textbf{Mat} --- l'arcane non numéroté, extérieur au pacte initial
(\og la cinquième puissance, celle qu'on ne doit pas nommer \fg{}).
Une \og sixième maison \fg{} apparaît ensuite : le complexe
militaro-industriel américain, lié au \og Roi en jaune \fg{}, le faux
Howard Hughes d'\emph{Arcane majeur}.

\section{La dérivation historique des enseignes}

La filiation réelle des enseignes (latines $\rightarrow$ germaniques
$\rightarrow$ françaises) est documentée :

\begin{table}[h]
\centering
\begin{tabular}{ p{0.20\textwidth} p{0.20\textwidth} p{0.20\textwidth} p{0.20\textwidth} p{0.20\textwidth} }
\toprule
\textbf{Enseigne~\newline latine} & \textbf{Intermédiaire~\newline germanique} & \textbf{Couleur~\newline française} & \textbf{Témoignage} \\
\midrule
Épée 	& Feuille 	& \pique{} Pique 		& anglais \emph{spades} $<$ \emph{espadas} \\
Denier 	& Grelot 	& \carreau{} Carreau 	& rond~\newline $\rightarrow$ losange \\
Bâton 	& Gland 	& \trefle{} Trèfle 		& anglais \emph{clubs} $<$ \emph{bastos} \\
Coupe 	& Cœur 		& \coeur{} Cœur 		& correspondance universelle \\
\bottomrule
\end{tabular}
\caption{Filiation historique des enseignes latines vers les couleurs françaises.}
\end{table}

Le schéma ci-dessous superpose les deux systèmes : en trait plein noir, la
filiation historique ; en pointillés rouges, l'attribution propre à la BD.

\begin{center}
\begin{tikzpicture}[>={Stealth[length=2.2mm]},every node/.style={font=\small}]
% Colonne gauche : enseignes latines (avec archonte)
\node[align=right] (coupe)  at (0, 0)    {Coupe (Dyo)};
\node[align=right] (epee)   at (0,-1.3)  {Épée (Aker)};
\node[align=right] (baton)  at (0,-2.6)  {Bâton / Lance (Reka)};
\node[align=right] (denier) at (0,-3.9)  {Denier (Erlin)};
% Colonne droite : couleurs françaises
\node (coeurL)   at (8.2, 0)    {\coeur{} Cœur};
\node (piqueL)   at (8.2,-1.3)  {\pique{} Pique};
\node (carreauL) at (8.2,-2.6)  {\carreau{} Carreau};
\node (trefleL)  at (8.2,-3.9)  {\trefle{} Trèfle};
% Filiation historique (noir, trait plein)
\draw[->] (coupe.east)  -- node[above,font=\scriptsize]{historique et BD} (coeurL.west);
\draw[->] (epee.east)   -- (piqueL.west);
\draw[->] (baton.east)  -- (trefleL.west);
\draw[->] (denier.east)  -- (carreauL.west);
% Attribution de la BD (rouge, pointillés)
\draw[->,red,dashed] (epee.east)   to[bend right=18] (carreauL.west);
\draw[->,red,dashed] (baton.east)  to[bend right=18] (piqueL.west);
\draw[->,red,dashed] (denier.east) to[bend right=18] (trefleL.west);
% Légende
\node[align=left,font=\scriptsize] at (4.1,-5.1)
  {\textcolor{black}{---} filiation historique \quad
   \textcolor{red}{- -} attribution BD};
\end{tikzpicture}
\end{center}

\section{Les discordances}

\subsection{L'inversion des couleurs (discordance principale)}

La BD respecte la logique étymologique uniquement pour les Coupes, et
permute les trois autres couleurs :

\begin{table}[h]
\centering
\begin{tabular}{p{0.30\textwidth} p{0.20\textwidth} p{0.20\textwidth} p{0.20\textwidth} p{0.20\textwidth}}
\toprule
\textbf{Enseigne} 		& \textbf{Couleur~\newline dans la BD} & \textbf{Couleur~\newline classique} & \textbf{Verdict} \\
\midrule
Coupes (Dyo) 			& \coeur{} Cœur 		& \coeur{} Cœur 		& \textcolor{green!55!black}{conforme} \\
Lances / Bâtons (Reka) 	& \pique{} Pique 		& \trefle{} Trèfle 		& \textcolor{red}{inversé} \\
Épées (Aker) 			& \carreau{} Carreau 	& \pique{} Pique 		& \textcolor{red}{inversé} \\
Deniers (Erlin) 		& \trefle{} Trèfle 		& \carreau{} Carreau 	& \textcolor{red}{inversé} \\
\bottomrule
\end{tabular}
\caption{Couleurs de la série vs correspondance classique (Rider-Waite, Papus, usage divinatoire).}
\end{table}

L'inversion la plus hétérodoxe est \textbf{Épées $\rightarrow$ Carreau} :
épées $\rightarrow$ piques est quasi universel (le mot \og pique \fg{}
désigne l'arme à pointe, par analogie avec l'épée latine ; l'anglais
\emph{spades} dérive de l'espagnol \emph{espadas}). À l'inverse,
\textbf{Lances $\rightarrow$ Pique} a une logique interne séduisante --- une
pique \emph{est} une arme d'hast, et le remplacement de \og Bâtons \fg{} par
\og Lances \fg{} dès le tome 1 semble fait pour autoriser ce jeu de mots.
Quant à \textbf{Deniers $\rightarrow$ Trèfle}, ce n'est pas une invention de
Pécau : la controverse occultiste est réelle (voir 4.2).

\subsection{La controverse occultiste Deniers / Bâtons}

Deux écoles s'affrontent depuis le XIX\ieme{} siècle sur l'attribution des
deux enseignes \og rondes \fg{} :
\begin{itemize}
\item \textbf{Papus} (et la majorité des auteurs) : Deniers $\rightarrow$
Carreaux, Bâtons $\rightarrow$ Trèfles ;
\item \textbf{Chevalier et Gheerbrant} (\emph{Dictionnaire des symboles}) :
Bâtons $\rightarrow$ Carreaux, Deniers $\rightarrow$ Trèfles.
\end{itemize}
Le choix de la série (Deniers $\rightarrow$ Trèfle) rejoint donc la seconde
école. Mais aucune tradition identifiée ne soutient
Épées $\rightarrow$ Carreau : sur ce point, la série crée sa propre
cartomancie.

\clearpage

\subsection{Hiérarchie des arcanes mineurs inversée}

Dans l'univers de Pécau, les agents (\og soldats \fg{}) sont classés en
force croissante : 2, 3, 4, 5, 6, 7, 8, 9, as, valet, cavalier, roi/reine
--- l'as vaut plus que le 9, et les figures plus encore. Dans le
\textbf{jeu de tarot réel}, la valeur est décroissante : le 21 est le plus
fort des atouts, et pour les enseignes c'est le roi qui domine. La BD
emprunte au tarot de Marseille son imagerie, mais sa hiérarchie est une
construction interne, proche des jeux de cartes modernes.

\subsection{Le Mat, l'Excuse et le Joker}

Le cas Guillaume de Lecce est cohérent avec la carte qu'il incarne. Le
\textbf{Mat}, arcane non numéroté du tarot de Marseille, devient
\textbf{l'Excuse} dans le jeu de tarot --- carte jouable à tout moment, hors
hiérarchie --- et son écho moderne est le \textbf{Joker} des jeux de
52 cartes, carte \og générique \fg{} échappant aux règles communes. Un
archonte extérieur au pacte des quatre ivoires, imprévisible et non
numéroté, est donc fidèle à la symbolique de sa carte. Né torturé à
Palerme, il est présenté comme la \og conséquence directe \fg{} de la
tentative d'Aker de créer le \og roi du monde \fg{} en la personne de
Frédéric II.

\subsection{Anecdote de vérification : le wiki qui a vacillé}

Un miroir d'une ancienne version de l'article Wikipédia indiquait
\og Aker (représente les épées / \emph{Pique} / l'Impératrice) \fg{} --- la
correspondance classique, pas celle de la BD. La version actuelle a été
réalignée sur \og Carreau \fg{}. Illustration involontaire de la force de
l'association mentale épées $\rightarrow$ piques, y compris chez les
contributeurs d'encyclopédies.

\section{Confirmé, interprétatif, non tranché}

\paragraph{Confirmé (sources concordantes).}
Les triplets enseigne / couleur / atout des cinq archontes ; le
remplacement de \og Bâtons \fg{} par \og Lances \fg{} dès le tome 1 ;
l'appellation \og bouclier \fg{} de l'ivoire d'Erlin chez le chaman ; la
5\ieme{} maison (moines rouges) et la 6\ieme{} (complexe militaro-industriel) ;
l'inversion des couleurs par rapport à la tradition.

\paragraph{Interprétatif.}
Le caractère \og volontaire \fg{} des inversions : les glissements de
vocabulaire (Bouclier/Denier, Lance/Bâton) servent visiblement l'ancrage
historique des époques, mais l'inversion des couleurs, elle, ne correspond
à aucune réalité historique documentée --- c'est un choix narratif assumé
(avec le jeu de mots Lance/Pique). De même, l'idée que les quatre atouts
\og incarnent l'ordre temporel et spirituel \fg{} face au chaos du Mat est
une lecture cohérente, mais aucun texte source ne la formule.

\paragraph{Non trouvé.}
Aucun billet de blog ni fil de forum de fans ne débattant explicitement de
l'inversion enseignes / couleurs. Une vérification primaire sur les albums
(scène de remise des ivoires au tome 1 ; pages de garde reprenant les
enseignes) resterait la source définitive.

\section{Sources}

\begin{itemize}
\item Wikipédia : \og L'Histoire secrète \fg{}, \og Arcane majeur \fg{},
\og Tarot Rider-Waite \fg{}, articles \og Pique \fg{}, \og Carreau \fg{},
\og Trèfle \fg{}, \og Excuse \fg{}, \og Joker \fg{} :
\href{https://fr.wikipedia.org/wiki/L'Histoire_secr%C3%A8te}{fr.wikipedia.org}
\item Miroir fr-academic (ancienne version de l'article) :
\href{https://fr-academic.com/dic.nsf/frwiki/932400}{fr-academic.com}
\item Babelio, critiques du tome 1 \emph{Genèse} (citation du discours de Père Loup) :
\href{https://www.babelio.com/livres/Pecau-LHistoire-Secrete-Tome-1--Genese/48758}{babelio.com}
\item Top Comics et BDThèque, fiches et avis concordants :
\href{https://topcomics.fr/l-histoire-secrete-jean-pierre-pecau-revisite-3-000-ans-d-histoire-avis-critique}{topcomics.fr},
\href{https://www.bdtheque.com/series/4490/l-histoire-secrete}{bdtheque.com}
\item Blog \emph{Les Secrets du Tarot}, \og Secrets des 56 arcanes mineurs \fg{}
(controverse Papus vs Chevalier \& Gheerbrant) :
\href{http://secretsdutarot.blogspot.com/2012/11/secrets-des-56-arcanes-mineurs.html}{secretsdutarot.blogspot.com}
\item Gaudry.be, \og L'Histoire secrète \fg{} (5\ieme{} et 6\ieme{} maisons) :
\href{https://www.gaudry.be/bd/l-histoire-secrete/}{gaudry.be}
\end{itemize}

\end{document}